% ===================== PREAMBULE COMMUN (à copier VERBATIM en tête de CHAQUE .tex) =====================
\documentclass[11pt,a4paper]{article}
\usepackage[utf8]{inputenc}
\usepackage[T1]{fontenc}
\usepackage{lmodern}
\usepackage[french]{babel}
\usepackage{amsmath,amssymb,mathtools}
\usepackage{icomma}
\usepackage{xcolor}
\usepackage{colortbl}
\usepackage{tikz}
\usepackage{pgfplots}
\usepackage[most]{tcolorbox}
\usepackage{enumitem}
\usepackage{array}
\usepackage{booktabs}
\usepackage{multicol}
\usepackage{fancyhdr}
\usepackage{caption}
\captionsetup{labelformat=empty,justification=centering,font=small}
\usepackage[hidelinks]{hyperref}
\usetikzlibrary{arrows.meta,positioning,calc,angles,quotes,patterns,backgrounds,shapes.geometric}
\pgfplotsset{compat=1.18}
\usepackage[a4paper,margin=2.2cm,headheight=26pt,headsep=10pt]{geometry}

\definecolor{primary}{RGB}{219,54,77}
\definecolor{primarydark}{RGB}{150,28,46}
\definecolor{primarylight}{RGB}{250,224,228}
\definecolor{defcol}{RGB}{0,114,178}
\definecolor{thmcol}{RGB}{0,140,120}
\definecolor{excol}{RGB}{0,135,90}
\definecolor{methcol}{RGB}{90,90,100}
\definecolor{courbe}{RGB}{40,90,200}
\definecolor{courbeB}{RGB}{210,60,40}
\definecolor{courbeC}{RGB}{30,150,80}
\definecolor{axiscol}{RGB}{60,60,60}
\definecolor{gridcol}{RGB}{200,205,215}

\tcbset{boxbase/.style={enhanced,breakable,boxrule=0.9pt,arc=2.5pt,
   left=3mm,right=3mm,top=2.3mm,bottom=2.3mm,
   fonttitle=\bfseries,coltitle=white,toptitle=0.6mm,bottomtitle=0.6mm}}
\newtcolorbox{methode}[1][]{boxbase,colback=methcol!7,colframe=methcol,sharp corners,
  title={\textsc{Méthode}\ifx\relax#1\relax\relax\else\textmd{ --- \itshape #1}\fi}}
\newtcolorbox{exemple}[1][]{boxbase,colback=excol!5,colframe=excol,
  title={\textsc{Exemple}\ifx\relax#1\relax\relax\else\textmd{ : \itshape #1}\fi}}
\newtcolorbox{idee}[1][]{boxbase,colback=defcol!6,colframe=defcol,
  title={\textsc{L'essentiel}\ifx\relax#1\relax\relax\else\textmd{ : \itshape #1}\fi}}
\newtcolorbox{piege}[1][]{boxbase,colback=primary!7,colframe=primary,
  title={\textsc{Piège}\ifx\relax#1\relax\relax\else\textmd{ : \itshape #1}\fi}}

\pgfplotsset{repere/.style={axis lines=middle,
    axis line style={-{Stealth[length=2.2mm]},axiscol,line width=0.8pt},
    label style={font=\footnotesize},tick label style={font=\scriptsize},
    every axis x label/.style={at={(current axis.right of origin)},anchor=north west},
    every axis y label/.style={at={(current axis.above origin)},anchor=south east},
    xlabel={$x$},ylabel={$y$},grid=both,minor tick num=0,
    major grid style={gridcol,line width=0.4pt},samples=200,clip=false}}

\pagestyle{fancy}\fancyhf{}
\fancyhead[L]{\small\textcolor{primarydark}{\textbf{Fiche 4} — Les fonctions}}
\fancyhead[R]{\small\textcolor{primarydark}{\textsc{Carnet d'entraînement}}}
\fancyfoot[C]{\small\thepage}
\renewcommand{\headrulewidth}{0.8pt}
\renewcommand{\headrule}{\hbox to\headwidth{\color{primary}\leaders\hrule height \headrulewidth\hfill}}
\usepackage{titlesec}
\titleformat{\section}{\Large\bfseries\color{primarydark}}{\thesection}{0.6em}{}
\titleformat{\subsection}{\large\bfseries\color{primary}}{\thesubsection}{0.6em}{}

\newcommand{\R}{\mathbb{R}}\newcommand{\Z}{\mathbb{Z}}\newcommand{\N}{\mathbb{N}}
\newcommand{\Df}{D_f}
\newcommand{\croit}{\textcolor{thmcol}{\Large$\nearrow$}}
\newcommand{\decroit}{\textcolor{thmcol}{\Large$\searrow$}}

\newcommand{\exercice}[1]{\medskip\noindent{\color{primarydark}\bfseries\large Exercice #1}\par\nopagebreak\smallskip}
\newcommand{\titrecarnet}[3]{%
\begin{center}\begin{tikzpicture}
\node[rectangle,rounded corners=7pt,fill=primary,text=white,minimum width=15cm,
  minimum height=1.7cm,align=center,inner sep=8pt]
  {{\LARGE\bfseries #1}\\[2pt]{\normalsize #2}};
\end{tikzpicture}\end{center}
\begin{center}\itshape\color{primarydark}#3\end{center}\medskip}

% ---- RÈGLES D'USAGE DES MACROS ----
% * \croit et \decroit s'emploient en mode TEXTE (dans une cellule de tableau),
%   JAMAIS entre $...$ (ils contiennent déjà un $\nearrow$).
% * \titrecarnet{Titre}{Sous-titre}{Ligne en italique} : bandeau de titre.
% * \exercice{1} : titre d'exercice.
% * Boîtes : \begin{idee}...\end{idee}, \begin{exemple}[titre]...\end{exemple},
%   \begin{methode}[titre]...\end{methode}, \begin{piege}[titre]...\end{piege}.
% Le corps du document commence après \begin{document}.
\begin{document}
\titrecarnet{Thème 5 — Trigonométrie, exponentielle \& logarithme}{Tutoriel}{Périodes, parités, formulaire trigonométrique, propriétés exp/log.}

\section*{1. Fonctions trigonométriques}

\begin{idee}[à connaître par cœur]
\begin{center}\renewcommand{\arraystretch}{1.2}
\begin{tabular}{lccc}
\toprule
 & période & parité & racines \\
\midrule
$\sin x$ & $2\pi$ & impaire & $x=k\pi$ \\
$\cos x$ & $2\pi$ & paire & $x=\tfrac{\pi}{2}+k\pi$ \\
$\tan x$ & $\pi$ & impaire & $x=k\pi$ \\
\bottomrule
\end{tabular}
\end{center}
Bornes : $-1\leq\sin x\leq1$, $-1\leq\cos x\leq1$. \ Réciproques : $\arcsin$ et $\arccos$ définies sur $[-1,1]$, $\arctan$ sur $\R$.
\end{idee}

\begin{tcolorbox}[boxbase,colback=primary!6,colframe=primary,title={\textsc{Formulaire trigonométrique}}]
\begin{multicols}{2}\small
\begin{itemize}[leftmargin=1.2em,topsep=0pt,itemsep=2pt]
  \item $\cos^2 x+\sin^2 x=1$
  \item $\sin(2x)=2\sin x\cos x$
  \item $\cos(2x)=\cos^2 x-\sin^2 x$
  \item $\sin(a+b)=\sin a\cos b+\cos a\sin b$
  \item $\cos(a+b)=\cos a\cos b-\sin a\sin b$
\end{itemize}
\end{multicols}
\end{tcolorbox}

Pour \emph{réduire un angle}, on retranche des tours entiers : $\cos(x+2k\pi)=\cos x$.

\section*{2. Exponentielle et logarithme}

$e^x$ est définie sur $\R$, toujours $>0$, croissante, $e^0=1$. Le logarithme est sa \textbf{réciproque} : $\ln x$ (défini sur $\,]0,+\infty[$), avec $\ln 1=0$, $\ln e=1$. Leurs courbes sont symétriques par rapport à la droite $y=x$.

\begin{tcolorbox}[boxbase,colframe=thmcol,colback=thmcol!6,title={\textsc{Propriétés indispensables}}]
\begin{multicols}{2}\small
\begin{itemize}[leftmargin=1.2em,topsep=0pt,itemsep=2pt]
  \item $e^a\cdot e^b=e^{a+b}$, \ $(e^a)^n=e^{na}$, \ $e^{\ln x}=x$
  \item $\ln(xy)=\ln x+\ln y$
  \item $\ln\dfrac{x}{y}=\ln x-\ln y$, \ $\ln x^n=n\ln x$
  \item Changement de base : $\log_a x=\dfrac{\ln x}{\ln a}$
\end{itemize}
\end{multicols}
\end{tcolorbox}

\begin{idee}[croissance exponentielle et temps de doublement]
Si une grandeur \emph{double} tous les $T$, alors après $n$ doublements elle est multipliée par $2^n$ (et il s'écoule $nT$). Modèle : $N(t)=N_0\cdot 2^{t/T}$. (Une \emph{demi-vie} suit $N(t)=N_0\cdot\big(\tfrac12\big)^{t/T}$.)
\end{idee}

\begin{center}
\begin{tikzpicture}
\begin{axis}[repere,width=8cm,height=4.8cm,xmin=-2.6,xmax=3.6,ymin=-2.6,ymax=3.6,
   xtick={-2,2},ytick={-2,2},samples=120,restrict y to domain=-2.6:3.6,title={\footnotesize $e^x$ et $\ln x$ symétriques / $y=x$}]
  \addplot[courbeB,line width=0.9pt,dashed,domain=-2.6:3.6]{x};
  \addplot[courbe,line width=1.2pt,domain=-2.5:1.25]{exp(x)};
  \addplot[courbeC,line width=1.2pt,domain=0.09:3.5]{ln(x)};
  \node[courbe,font=\scriptsize]at(axis cs:-1.5,2.6){$e^x$};
  \node[courbeC,font=\scriptsize]at(axis cs:2.7,1.3){$\ln x$};
\end{axis}
\end{tikzpicture}
\end{center}
\end{document}
